\documentclass[11pt, ukrainian]{article}
%\setlength\textwidth{190mm} \setlength\textheight{237mm}%\setlength\voffset{-18mm} \setlength\hoffset{-33mm}\usepackage[T2A]{fontenc}
\usepackage[utf8]{inputenc}
\usepackage[ukrainian]{babel}
\usepackage{graphicx}
\usepackage[Export]{adjustbox}
\adjustboxset{max size={1.0\linewidth}{1.0\paperheight}}
\graphicspath{{data/pages/}} 
\usepackage{geometry}
\geometry{top=1.5cm,bottom=1.5cm,left=1.2cm,right=1.2cm}
\pagestyle{empty}\setlength\parindent{0mm}
\begin{document}
{\Large \center  Варіанти завдань до програмного проєкту №3\par Системний аналіз, 2024/25 н.р.\par \bigskip\bigskip}

\vspace{2mm}
У всіх варіантах діагоналі та інші лінії, що обмежують частину матриці, яку треба обійти, ВХОДЯТЬ до області, що обходиться. Наприклад, трикутник над головною діагоналлю включає саму цю головну діагональ.

\vspace{2mm}

Вважати, що лівий верхній кут матриці має координати~$(0,0)$.

\vspace{2mm}

Якщо попередні вимоги на початок обходу залишають неоднозначність, то розглядати ті, з яких можливий обхід усієї зазначеної області, якщо і після цього залишається неоднозначність, то перевагу віддавати тому елементу матриці з наявних альтернатив, що має лексикографічно менші індекси.

\vspace{2mm}

Якщо перший крок обходу можна зробити кількома способами, то розглядати ті, за яких можливий обхід усієї зазначеної області, якщо і після цього залишається неоднозначність, то перевагу на першому кроці віддавати переходу до того елемента матриці з наявних альтернатив, що має лексикографічно менші індекси.

\vspace{2mm}

%begin
{\Large \center  Варіант  60\par \bigskip\bigskip}

За квадратною матрицею, заданою об'єктом типу numpy.ndarray, надрукувати послідовність її елементів обходом верхнього чверть-трикутника матриці ``змійкою'' від центрального кута діагоналями (паралельно головній діагоналі).
%z03 Z p09 n++


(Перший кро змійка робить по діагоналі вліво вгору.)


%end
%begin
{\Large \center  Варіант  61\par \bigskip\bigskip}

За квадратною матрицею, заданою об'єктом типу numpy.ndarray, надрукувати послідовність її елементів обходом лівого чверть-трикутника матриці ``змійкою'' від верхнього лівого кута матриці діагоналями (перпендикулярно головній діагоналі).
%z06 Z p07 n+-


(Перший крок змійка робить вертикально вниз.)


%end
%begin
{\Large \center  Варіант  62\par \bigskip\bigskip}

За квадратною матрицею, заданою об'єктом типу numpy.ndarray, надрукувати послідовність її елементів обходом трикутника над побічною діагоналлю матриці спіраллю від верхнього лівого кута матриці за годинниковою стрілкою.
%z07 S p07 n-





%end
%begin
{\Large \center  Варіант  63\par \bigskip\bigskip}

За квадратною матрицею, заданою об'єктом типу numpy.ndarray, надрукувати послідовність її елементів обходом центрального ромба матриці спіраллю від його правого кута проти годинникової стрілки.
%z00 S p01 n+





%end
%begin
{\Large \center  Варіант  64\par \bigskip\bigskip}

За квадратною матрицею, заданою об'єктом типу numpy.ndarray, надрукувати послідовність її елементів обходом центрального ромба матриці ``змійкою'' від його верхнього кута горизонталями.
%z00 Z p04 n0+


(Вважати, що перша горизонталь йде зліва направо.)


%end
%begin
{\Large \center  Варіант  65\par \bigskip\bigskip}

За квадратною матрицею, заданою об'єктом типу numpy.ndarray, надрукувати послідовність її елементів обходом правого чверть-трикутника матриці ``змійкою'' від нижнього правого кута матриці вертикалями.
%z04 Z p06 n+0





%end
%begin
{\Large \center  Варіант  66\par \bigskip\bigskip}

За квадратною матрицею, заданою об'єктом типу numpy.ndarray, надрукувати послідовність її елементів обходом нижнього чверть-трикутника матриці спіраллю від нижнього лівого кута матриці за годинниковою стрілкою.
%z05 S p05 n-





%end
%begin
{\Large \center  Варіант  67\par \bigskip\bigskip}

За квадратною матрицею, заданою об'єктом типу numpy.ndarray, надрукувати послідовність її елементів обходом нижнього чверть-трикутника матриці спіраллю від нижнього правого кута матриці проти годинникової стрілки.
%z05 S p06 n+





%end
%begin
{\Large \center  Варіант  68\par \bigskip\bigskip}

За квадратною матрицею, заданою об'єктом типу numpy.ndarray, надрукувати послідовність її елементів обходом нижнього чверть-трикутника матриці ``змійкою'' від центрального кута діагоналями (перпендикулярно головній діагоналі).
%z05 Z p09 n+-





%end
%begin
{\Large \center  Варіант  69\par \bigskip\bigskip}

За квадратною матрицею, заданою об'єктом типу numpy.ndarray, надрукувати послідовність її елементів обходом трикутника над побічною діагоналлю матриці ``змійкою'' від верхнього правого кута матриці вертикалями.
%z07 Z p08 n+0


(Перший крок змійка робить ліворуч.)

%end
%begin
{\Large \center  Варіант  70\par \bigskip\bigskip}

За квадратною матрицею, заданою об'єктом типу numpy.ndarray, надрукувати послідовність її елементів обходом трикутника над головною діагоналлю матриці ``змійкою'' від верхнього правого кута матриці діагоналями (паралельно головній діагоналі).
%z01 Z p08 n++





%end
%begin
{\Large \center  Варіант  71\par \bigskip\bigskip}

За квадратною матрицею, заданою об'єктом типу numpy.ndarray, надрукувати послідовність її елементів обходом правого чверть-трикутника матриці ``змійкою'' від центрального кута вертикалями.
%z04 Z p09 n+0


(Вважати, що перша вертикаль йде вгору.)


%end
%begin
{\Large \center  Варіант  72\par \bigskip\bigskip}

За квадратною матрицею, заданою об'єктом типу numpy.ndarray, надрукувати послідовність її елементів обходом трикутника над побічною діагоналлю матриці ``змійкою'' від нижнього лівого кута матриці вертикалями.
%z07 Z p05 n+0





%end
%begin
{\Large \center  Варіант  73\par \bigskip\bigskip}

За квадратною матрицею, заданою об'єктом типу numpy.ndarray, надрукувати послідовність її елементів обходом трикутника під головною діагоналлю матриці спіраллю від нижнього лівого кута матриці за годинниковою стрілкою.
%z02 S p05 n-





%end
%begin
{\Large \center  Варіант  74\par \bigskip\bigskip}

За квадратною матрицею, заданою об'єктом типу numpy.ndarray, надрукувати послідовність її елементів обходом трикутника під головною діагоналлю матриці ``змійкою'' від нижнього правого кута матриці діагоналями (паралельно головній діагоналі).
%z02 Z p06 n++





%end
%begin
{\Large \center  Варіант  75\par \bigskip\bigskip}

За квадратною матрицею, заданою об'єктом типу numpy.ndarray, надрукувати послідовність її елементів обходом правого чверть-трикутника матриці ``змійкою'' від верхнього правого кута матриці горизонталями.
%z04 Z p08 n0+


(Перший крок змійка робить по діагоналі вліво та вниз.)


%end
%begin
{\Large \center  Варіант  76\par \bigskip\bigskip}

За квадратною матрицею, заданою об'єктом типу numpy.ndarray, надрукувати послідовність її елементів обходом трикутника під побічною діагоналлю матриці спіраллю від нижнього правого кута матриці проти годинникової стрілки.
%z08 S p06 n+





%end
%begin
{\Large \center  Варіант  77\par \bigskip\bigskip}

За квадратною матрицею, заданою об'єктом типу numpy.ndarray, надрукувати послідовність її елементів обходом трикутника під побічною діагоналлю матриці ``змійкою'' від верхнього правого кута матриці діагоналями (перпендикулярно головній діагоналі).
%z08 Z p08 n+-

(Перший крок змійка робить вертикально вгору.)



%end
%begin
{\Large \center  Варіант  78\par \bigskip\bigskip}

За квадратною матрицею, заданою об'єктом типу numpy.ndarray, надрукувати послідовність її елементів обходом трикутника над головною діагоналлю матриці спіраллю від нижнього правого кута матриці за годинниковою стрілкою.
%z01 S p06 n-





%end
%begin
{\Large \center  Варіант  79\par \bigskip\bigskip}

За квадратною матрицею, заданою об'єктом типу numpy.ndarray, надрукувати послідовність її елементів обходом верхнього чверть-трикутника матриці ``змійкою'' від верхнього лівого кута матриці вертикалями.
%z03 Z p07 n+0


(Перший крок змійка робить горизонтально вправо.)


%end
%begin
{\Large \center  Варіант  80\par \bigskip\bigskip}

За квадратною матрицею, заданою об'єктом типу numpy.ndarray, надрукувати послідовність її елементів обходом трикутника під головною діагоналлю матриці спіраллю від верхнього лівого кута матриці проти годинникової стрілки.
%z02 S p07 n+





%end
%begin
{\Large \center  Варіант  81\par \bigskip\bigskip}

За квадратною матрицею, заданою об'єктом типу numpy.ndarray, надрукувати послідовність її елементів обходом лівого чверть-трикутника матриці ``змійкою'' від центрального кута вертикалями.
%z06 Z p09 n+0


(Вважати, що перша вертикаль йде вниз.)


%end
%begin
{\Large \center  Варіант  82\par \bigskip\bigskip}

За квадратною матрицею, заданою об'єктом типу numpy.ndarray, надрукувати послідовність її елементів обходом трикутника під побічною діагоналлю матриці ``змійкою'' від нижнього лівого кута матриці діагоналями (паралельно головній діагоналі).
%z08 Z p05 n++





%end
%begin
{\Large \center  Варіант  83\par \bigskip\bigskip}

За квадратною матрицею, заданою об'єктом типу numpy.ndarray, надрукувати послідовність її елементів обходом верхнього чверть-трикутника матриці ``змійкою'' від верхнього правого кута матриці діагоналями (перпендикулярно головній діагоналі).
%z03 Z p08 n+-





%end
%begin
{\Large \center  Варіант  84\par \bigskip\bigskip}

За квадратною матрицею, заданою об'єктом типу numpy.ndarray, надрукувати послідовність її елементів обходом лівого чверть-трикутника матриці ``змійкою'' від нижнього лівого кута матриці діагоналями (паралельно головній діагоналі).
%z06 Z p05 n++


(Перший крок змійка робить вертикально вгору.)


%end
%begin
{\Large \center  Варіант  85\par \bigskip\bigskip}

За квадратною матрицею, заданою об'єктом типу numpy.ndarray, надрукувати послідовність її елементів обходом центрального ромба матриці ``змійкою'' від його лівого кута вертикалями.
%z00 Z p02 n+0

(Вважати, що перша вертикаль йде вниз.)



%end
%begin
{\Large \center  Варіант  86\par \bigskip\bigskip}

За квадратною матрицею, заданою об'єктом типу numpy.ndarray, надрукувати послідовність її елементів обходом трикутника над головною діагоналлю матриці ``змійкою'' від верхнього лівого кута матриці діагоналями (перпендикулярно головній діагоналі).
%z01 Z p07 n+-

(Перший крок змійка робить горизонтально.)



%end
%begin
{\Large \center  Варіант  87\par \bigskip\bigskip}

За квадратною матрицею, заданою об'єктом типу numpy.ndarray, надрукувати послідовність її елементів обходом центрального ромба матриці ``змійкою'' від його нижнього кута горизонталями.
%z00 Z p03 n0+

(Перша горизонталь йде зліва направо.)



%end
%begin
{\Large \center  Варіант  88\par \bigskip\bigskip}

За квадратною матрицею, заданою об'єктом типу numpy.ndarray, надрукувати послідовність її елементів обходом трикутника під побічною діагоналлю матриці спіраллю від нижнього лівого кута матриці за годинниковою стрілкою.
%z08 S p05 n-





%end
%begin
{\Large \center  Варіант  89\par \bigskip\bigskip}

За квадратною матрицею, заданою об'єктом типу numpy.ndarray, надрукувати послідовність її елементів обходом трикутника під побічною діагоналлю матриці спіраллю від верхнього правого кута матриці проти годинникової стрілки.
%z08 S p08 n+





%end
%begin
{\Large \center  Варіант  90\par \bigskip\bigskip}

За квадратною матрицею, заданою об'єктом типу numpy.ndarray, надрукувати послідовність її елементів обходом верхнього чверть-трикутника матриці ``змійкою'' від верхнього правого кута матриці горизонталями.
%z03 Z p08 n0+





%end
%begin
{\Large \center  Варіант  91\par \bigskip\bigskip}

За квадратною матрицею, заданою об'єктом типу numpy.ndarray, надрукувати послідовність її елементів обходом нижнього чверть-трикутника матриці спіраллю від нижнього правого кута матриці проти годинникової стрілки.
%z05 S p06 n+





%end
%begin
{\Large \center  Варіант  92\par \bigskip\bigskip}

За квадратною матрицею, заданою об'єктом типу numpy.ndarray, надрукувати послідовність її елементів обходом трикутника під головною діагоналлю матриці ``змійкою'' від верхнього лівого кута матриці діагоналями (паралельно головній діагоналі).
%z02 Z p07 n++




%end
%begin
{\Large \center  Варіант  93\par \bigskip\bigskip}

За квадратною матрицею, заданою об'єктом типу numpy.ndarray, надрукувати послідовність її елементів обходом трикутника над побічною діагоналлю матриці ``змійкою'' від нижнього лівого кута матриці вертикалями.
%z07 Z p05 n+0



(Перший крок змійка робить по діагоналі.)


%end
%begin
{\Large \center  Варіант  94\par \bigskip\bigskip}

За квадратною матрицею, заданою об'єктом типу numpy.ndarray, надрукувати послідовність її елементів обходом правого чверть-трикутника матриці спіраллю від верхнього правого кута матриці за годинниковою стрілкою.
%z04 S p08 n-





%end
%begin
{\Large \center  Варіант  95\par \bigskip\bigskip}

За квадратною матрицею, заданою об'єктом типу numpy.ndarray, надрукувати послідовність її елементів обходом трикутника під побічною діагоналлю матриці ``змійкою'' від нижнього правого кута матриці діагоналями (перпендикулярно головній діагоналі).
%z08 Z p06 n+-


(Перший крок змійка робить вертикально.)



%end
%begin
{\Large \center  Варіант  96\par \bigskip\bigskip}

За квадратною матрицею, заданою об'єктом типу numpy.ndarray, надрукувати послідовність її елементів обходом правого чверть-трикутника матриці ``змійкою'' від нижнього правого кута матриці вертикалями.
%z04 Z p06 n+0





%end
%begin
{\Large \center  Варіант  97\par \bigskip\bigskip}

За квадратною матрицею, заданою об'єктом типу numpy.ndarray, надрукувати послідовність її елементів обходом центрального ромба матриці ``змійкою'' від його правого кута діагоналями (паралельно головній діагоналі).
%z00 Z p01 n++





%end
%begin
{\Large \center  Варіант  98\par \bigskip\bigskip}

За квадратною матрицею, заданою об'єктом типу numpy.ndarray, надрукувати послідовність її елементів обходом трикутника над головною діагоналлю матриці ``змійкою'' від нижнього правого кута матриці діагоналями (перпендикулярно головній діагоналі).
%z01 Z p06 n+-


(Перший крок змійка робить по діагоналі вліво вгору.)


%end
%begin
{\Large \center  Варіант  99\par \bigskip\bigskip}

За квадратною матрицею, заданою об'єктом типу numpy.ndarray, надрукувати послідовність її елементів обходом трикутника під головною діагоналлю матриці спіраллю від нижнього лівого кута матриці проти годинникової стрілки.
%z02 S p05 n+





%end
\end{document}
